\chapter{Updating the bridge}
\label{ch:update}

An update is one gesture: a new image on the same volume. \prog keeps everything it needs in
the state directory --- \code{config.toml} and \code{bdseq.toml} --- and the container itself
holds nothing you would miss. What follows is the procedure, the point you can go back to, and
what to look at afterwards to know it worked.

\section{The procedure}

\code{docker-compose.yml} pins an exact image, deliberately: \code{latest} would swap the
bridge at whatever hour a pull happened to run. Updating is therefore \emph{editing that line},
not merely pulling.

\begin{lstlisting}
# 1. Note where you are. This is the rollback point.
docker compose config | grep image:

# 2. Edit docker-compose.yml: image: docker.io/gcorbaz/smartme_mqtt:<new tag>

# 3. Pull and restart. `up -d` recreates the container; the volume is untouched.
docker compose pull
docker compose up -d

# 4. Watch it come back.
docker compose logs -f --tail=50
\end{lstlisting}

\begin{note}
\textbf{The volume is what carries your configuration across the update}, and nothing else. A
new container reads the same \code{config.toml}, finds the mapping already confirmed, and
starts publishing without asking you anything. That is asserted against the real image on every
CI run: a container is configured through its own screens, removed, and a \emph{second}
container is started on the same volume --- it must come up publishing and must not rewrite the
file.
\end{note}

\begin{note}
\textbf{If the update also changes what the bridge publishes, do both in the same stop.} An image
that carries a new contract and a rename of the published identity are two crossings of the same
window, and a host files history under the names it receives: each crossing made separately breaks
the series once. Stop the container, edit the compose file \emph{and} the configuration --- and the
state files the identity keys, see section~\ref{sec:changecost} --- then start once. Setting
\code{mapping\_confirmed} back to \code{false} while you are in the file is worth the extra click:
the bridge then comes up serving its screens and publishing nothing, so you read the mapping before
a host discovers it rather than after.
\end{note}

\section{The rollback point, and its one trap}

The rollback point is \textbf{the image tag you were running}, which step~1 above writes down.
Going back is the same gesture with the old tag:

\begin{lstlisting}
# Put the previous tag back in docker-compose.yml, then:
docker compose pull
docker compose up -d
\end{lstlisting}

\begin{warning}
\textbf{A configuration written by a newer \prog is refused by an older one, and this is the
trap to know before you need it.} \code{config.toml} carries a \code{schema\_version}. A newer
build may add keys; an older build refuses a file whose version it does not recognise, rather
than reading it by guesswork and starting on settings nobody wrote. So a rollback \emph{across
a schema change} leaves you with a bridge that starts, serves its web UI, and refuses to
publish until the configuration is re-entered.

\textbf{What to do about it:} before an update that changes the schema, copy the state
directory. \code{cp -a ./data ./data.bak} costs nothing and is the only thing that makes the
rollback complete rather than partial.

The schema version is printed in the file itself, and \S\ref{sec:changecost} lists which
versions this build reads.
\end{warning}

\section{Verifying the update}

Three signals, in the order they become available. \textbf{Each says something different}, and
the first two can be true while the third is not.

\begin{center}
\begin{tabular}{@{}>{\bfseries}l p{8.6cm}@{}}
\toprule
\normalfont\textbf{Signal} & \textbf{What it proves, and what it does not} \\
\midrule
\code{docker ps} & the container is up and its health is \emph{healthy}. This proves the
process runs and its poll loop is ticking --- \emph{not} that anything reaches your broker \\
\code{/healthz} & \code{"status":"running"}, plus \code{sink\_connected}, the failed sources,
the degraded meters and the dropped readings. This is the machine-facing answer, and the
sink line is the one the health status deliberately ignores \\
The state screen & the version and contract at the foot of the page --- \textbf{this is where
you confirm the new image is the one running} --- and the per-meter table showing fresh
publications \\
\bottomrule
\end{tabular}
\end{center}

\begin{note}
\textbf{An unreachable broker will not make the container unhealthy, and that is on purpose.}
The health status turns unhealthy for exactly one fault: a publishing bridge whose poll loop
has stopped ticking, which a restart repairs. A broker outage, a refused credential and a
frozen meter are all reported in the body of \code{/healthz} and on the screens, and none of
them restarts the container --- restarting would destroy every meter's Sparkplug session and
repair nothing.
\end{note}

\begin{note}
\textbf{The one thing to check on the SCADA side} is that each device came back with a
\code{DBIRTH} and is publishing again. \prog opens a new session on every start, so the host
sees a death and a birth; a host that shows a device still dead a minute after the update has
not received the birth, and the broker is where to look.
\end{note}
